\documentclass[11pt]{article}

\usepackage[T1]{fontenc}
\usepackage[utf8]{inputenc}
\usepackage[margin=1in]{geometry}
\usepackage{amsmath}
\usepackage{booktabs}

\title{Supplementary Representative Counts for Table 4}
\date{}

\begin{document}

\maketitle

This note records the representative counts for the three known algebraic
realizations listed in Table~4 of the main paper.

\begin{table}[ht]
\centering
\begin{tabular}{@{}llr@{}}
\toprule
Code & Source & Representatives \\
\midrule
$\mathsf N_{\mathrm I}^{+}(14;5,2)$ & Orevkov and Fiedler, $k=4$ & 2 \\
$\mathsf N_{\mathrm I}^{+}(14;1,0)$ & Orevkov, $k=6$ & 1 \\
$\mathsf N_{\mathrm I}^{+}(12;0,0)$ & Brugall\'e & 1 \\
\bottomrule
\end{tabular}
\caption{Representative counts for known algebraic realizations outside the primitive patchworking witness list.}
\end{table}

\end{document}